\section{Program Transformation}
\label{sec:transformation}

We now present program transformation techniques that replace libc streams in
C2Rust-generated code with Rust std streams using the analysis results.
%
The transformation consists of four steps:
%
transforming types (\cref{sec:transformation:types}),
%
transforming stream construction expressions
(\cref{sec:transformation:streams}),
%
transforming assignments and function calls
(\cref{sec:transformation:assignments}),
%
and propagating errors (\cref{sec:transformation:errors}).
%
% While we apply all these changes in a single pass by traversing the program
% once, we explain them as separate steps for clarity.

\subsection{Transforming Types}
\label{sec:transformation:types}

In this step, we replace each type annotation of \C{*mut\ FILE} (Rust syntax for
\C{FILE*}) with an appropriate type from Rust std.
%
We define the stream types provided by Rust std as follows:
%
\[
  %
  \small
  %
  \begin{array}{@{}r@{~}r@{~}l@{}}
  %
    B &::=& \tstdin\ |\ \tstdout\ |\ \tstderr\ |\ \tfile \\
  %
    &|& \tbufreader\ |\ \tbufwriter\ |\ \tchild \\
  %
    T &::=& \tread\ |\ \tbufread\ |\ \twrite\ |\ \tseek \\
  %
    \beta &::=& \overline{T} \\
  %
    \tau &::=& B\ |\ \tptr{B}\ |\ \tbox{\beta}\ |\ \tdynptr{\beta}\ |\
    \timpl{\beta} \\
  %
  \end{array}
  %
\]

$B$ is a \emph{base type}, where {\small$\tstdin$}, {\small$\tstdout$},
{\small$\tstderr$}, {\small$\tfile$}, and {\small$\tchild$} correspond to each
origin, and {\small$\tbufreader$} and {\small$\tbufwriter$} wrap files with
buffers.
%
We use \C{Child} for pipes instead of \C{ChildStdin} or \C{ChildStdout}, due to
the behavior of \C{pclose} in libc, which waits for the subprocess to terminate.
%
In Rust std, a subprocess is represented by \C{Child}, which can be waited on
and contains \C{ChildStdin} and \C{ChildStdout} as fields for I/O.
%
To ensure correct translation of \C{pclose}, we use \C{Child} as a stream
originating from a pipe.

$T$ is a trait corresponding to each capability except $\cclose$.
%
$\beta$ is a list of traits, representing bounds.

A type $\tau$ is either a base type, a raw pointer to a base type, \C{dyn}
wrapped by \C{Box}, \C{dyn} wrapped by a raw pointer, or \C{impl}.
%
% When $\beta$ after \C{dyn} contains multiple bounds, we define a subtrait
% inheriting all the bounds and use it as the bound.
%
% For \C{impl}, the bounds in $\beta$ are simply enumerated using \C{+}.
%
We use raw pointers instead of references because references require their
lifetimes to be decided, which is unrelated to I/O and outside the scope of this
work.
%
Existing techniques~\cite{emre2021translating, emre2023aliasing,
zhang2023ownership} can be applied to replace raw pointers with references.

\begin{algorithm}[t]
  \caption{Type decision}
  \label{alg:type}
  \setstretch{0.8}
  \footnotesize
  \Fn{decideType($l$)}{
    \lIf{$l$ is a parameter}{
      \Return{$\timpl{\boundof{\capset{l}}}$}\label{line:type2}
    }\ElseIf{$\#(\oriset{l})=1$}{\label{line:type3}
      \lIf{$\oriset{l}=\{\ostdin\}$}{
        $B:=\tstdin$\label{line:type4}
      }\lElseIf{$\oriset{l}=\{\ostdout\}$}{
        $B:=\tstdout$\label{line:type5}
      }\lElseIf{$\oriset{l}=\{\ostderr\}$}{
        $B:=\tstderr$\label{line:type6}
      }\ElseIf{$\oriset{l}=\{\ofile\}$}{\label{line:type7}
        \lIf{$\cwrite\not\in\capset{l}$}{
          $B:=\tbufreader$\label{line:type8}
        }\ElseIf{$\{\cread,\cbufread\}\cap\capset{l}=\emptyset$}{\label{line:type9}
          $B:=\tbufwriter$\;\label{line:type10}
        }\lElse{
          $B:=\tfile$\label{line:type11}
        }
      }\lElse{
        $B:=\tchild$\label{line:type12}
      }
      \lIf{$\cclose\in\capset{l}$}{
        \Return{$B$}\label{line:type13}
      }\lElse{
        \Return{$\tptr{B}$}\label{line:type14}
      }
    }\Else{\label{line:type15}
      \lIf{$\cclose\in\capset{l}$}{
        \Return{$\tbox{\boundof{\capset{l}}}$}\label{line:type16}
      }\lElse{
        \Return{$\tdynptr{\boundof{\capset{l}}}$}\label{line:type17}
      }
    }
  }
\end{algorithm}

\cref{alg:type} shows how we decide the type of each location $l$ based on its
origins and capabilities.
%
The type determined by this algorithm replaces the type annotation of the
location.
%
We treat the cases where $l$ is a parameter and where it is not differently.
%
We first discuss the non-parameter cases, which are further divided into
single-origin and multi-origin cases.

When only one origin is possible, we determine the base type according to the
origin (lines~\ref{line:type4}--\ref{line:type12}).
%
A notable case is when the origin is $\ofile$, where we also examine the
capabilities.
%
If the location is not used for writing, we use {\small$\tbufreader$}, and if it
is not used for reading, we use {\small$\tbufwriter$}, making the stream
buffered.
%
However, if the location is used for both reading and writing, we use
{\small$\tfile$} to support both kinds of operations.
%
Finally, if the location is used for closing, we assign the base type as the
type of the location (line~\ref{line:type13});
%
otherwise, we use a raw pointer (line~\ref{line:type14}).

When multiple origins are possible, we use \C{dyn} to allow storing different
base types in the location.
%
If the location is used for closing, we wrap the \C{dyn} type in a \C{Box}
(line~\ref{line:type16});
%
otherwise, we use a raw pointer (line~\ref{line:type17}).
%
Here, ${\it bound}$ converts the given set of capabilities into bounds by
mapping each capability to its corresponding trait, while ignoring $\cclose$.

Finally, when $l$ is a parameter, we use \C{impl} regardless of the number of
origins to give the most general type to the parameter (line~\ref{line:type2}).
%
We also do not check whether the stream is used for closing.
%
This is because we can handle this at the call site instead---passing an owning
value to close the stream, or a borrowing value to leave it open.
%
% Note that for each of \C{Read}, \C{BufRead}, \C{Write}, and \C{Seek}, if a type
% $A$ implements a trait, then \C{\&mut} $A$ also implements the same trait,
% allowing functions to accept both owning and borrowing values.

% For example, consider a function that reads the given stream. Even if the
% current version of the program always call this function with a file, it may
% need to call it with standard input in the future. If we use $\tfile$ as the
% parameter type, such change requires fixing the function signature, while using
% \C{impl} already allows passing standard input. As \C{impl} does not incur any
% run-time overhead, we have no reason to give up this flexibility.

In fact, since libc streams can be null pointers while Rust std streams cannot,
we need to use \C{Option}\CT{<}$\tau$\CT{>} instead of $\tau$ as the type of
each location.
%
An \C{Option} value is either \C{Some(}$v$\C{)}, representing a non-null stream,
or \C{None}, representing a null stream.
%
However, this issue of nullability arises not only in the context of I/O, but
also when handling any raw pointer;
%
previous studies that replace raw pointers with references also wrap references
in \C{Option}~\cite{emre2021translating, emre2023aliasing, zhang2023ownership}.
%
Thus, the remaining discussion omits \C{Option} to focus on I/O-specific
aspects.

\subsection{Transforming Stream Construction Expressions}
\label{sec:transformation:streams}

In this step, we transform each expression that constructs a libc stream into an
expression that constructs a Rust std stream.
%
We replace \C{stdin}, \C{stdout}, and \C{stderr} with \C{stdin()}, \C{stdout()},
and \C{stderr()} because standard streams are global variables in libc but are
returned from functions in Rust std.
%
We also replace \C{fopen} calls with calls to functions that return \C{File},
such as \C{File{::}open}, and replace \C{popen} calls with
\C{Command{::}new(..).spawn()}, which returns \C{Child}.

\subsection{Transforming Assignments and Function Calls}
\label{sec:transformation:assignments}

In this step, we transform assignments and function calls involving streams.
%
This is necessary because the types of the left value (or parameter) and the
right value (or argument) may differ.
%
For example, the left value can have more origins than the right value according
to the analysis, resulting in the left value having \C{Box} while the right
value has a base type.
%
Therefore, we must transform each right value to make it conform to the type of
the left value.

\begin{table}[t]
  \caption{Transformation of right values according to types}
  \label{tab:rvalue}
  \vspace{-1em}
  \centering
  $
  \footnotesize
  \begin{array}{@{~}c@{~}||@{~}c@{~}|@{~}c@{~}|@{~}c@{~}|@{~}c@{~}}
    & B_r & \tptr{B_r} & \tbox{\beta_r} & \tdynptr{\beta_r} \\ \hline\hline
    B_l & e\ ^\dagger & & & \\ \hline
    \tptr{B_l} & \eptr{e} & e & & \\ \hline
    \tbox{\beta_l} & {\tt Box{::}new}(e) & & e & \\ \hline
    \tdynptr{\beta_l} & \eptr{e} & e & e.{\tt as\_mut\_ptr}() & e \\
  \end{array}
  $
  \vspace{0.25em}

  {\scriptsize ${}^\dagger$
  ${\tt BufReader{::}new}(e)$ when $B_l=\tbufreader$ and $B_r=\tfile$;\\[-0.5em]
  ${\tt BufWriter{::}new}(e)$ when $B_l=\tbufwriter$ and $B_r=\tfile$}
  \vspace{-1em}
\end{table}

We first discuss the transformation of the right values of assignments, as
summarized in \cref{tab:rvalue}.
%
Rows represent the types of left values, and columns represent the types of
right values.
%
Each cell shows how the right value is transformed, where $e$ denotes the
original expression.
%
Since \C{impl} types only appear as parameter types when transforming function
calls, they are not considered in assignment transformations.

An empty cell in the table indicates that such a case never occurs.
%
Since the right value has more capabilities than the left, we never convert a
borrowing value to an owning value.
%
Likewise, as the right value has fewer origins, we never convert a \C{dyn} type
to a base type.

During the transformation, we only compare the categories of the types (i.e.,
base type, raw pointer, \C{Box\ dyn}, and \C{dyn} pointer) and do not check
whether the inner base types ($B_l$ and $B_r$) or bounds ($\beta_l$ and
$\beta_r$) are compatible.
%
This is because types that make conversion infeasible never occur, as types are
determined from the origins and capabilities computed by the analysis.
%
For example, we never transform {\small$\tstdin$} to {\small$\tstdout$}.
%
The only exception is when converting {\small$\tfile$} to {\small$\tbufreader$}
or {\small$\tbufwriter$}, where we use \C{BufReader{::}new(}$e$\C{)} or
\C{BufWriter{::}new(}$e$\C{)}, instead of $e$.

\begin{table}[t]
  \caption{Transformation of arguments according to types}
  \label{tab:arg}
  \vspace{-1em}
  \centering
  $\footnotesize
  \begin{array}{@{~}c@{~}|@{~}c@{~}|@{~}c@{~}|@{~}c@{~}}
    B & \tptr{B} & \tbox{\beta} & \tdynptr{\beta} \\ \hline\hline
    {\tt\&mut\ }e & e & e.{\tt as\_mut}() & e \\
  \end{array}$
  \vspace{-1.5em}
\end{table}

We now describe the transformation of arguments for user-defined function calls.
%
If the parameter of the function requires $\cclose$, then the argument remains
unchanged, as its type is guaranteed to be a base type or \C{Box}, and passing
the owning value allows the function to close the stream.
%
On the other hand, if the parameter does not require $\cclose$, we transform the
argument based on its type, as shown in \cref{tab:arg}.
%
For owning values, we take their addresses to prevent the stream from being
closed;
%
for borrowing values, we pass them as-is.

Finally, we discuss the transformation of libc API function calls.
%
The only two functions that close a stream are \C{fclose} and \C{pclose}.
%
We transform these into calls to \C{drop}, keeping the arguments unchanged.
%
Other function calls are translated into their Rust std equivalents.
%
For example, \C{fputc(}$v, e$\C{)} becomes $e'$\C{.write\_all(\&[}$v$\C{])}, where
$e'$ is the transformed form of $e$ according to \cref{tab:arg}.
%
The supplementary material~\cite{supp} lists all libc I/O functions and their
translations in Rust std.

Below is an example showing the result of the transformation up to this step:

\begin{lstlisting}[language=Rust]
// Before transformation
let x: *mut FILE = fopen(..);
let y: *mut FILE = popen(..);
let z: *mut FILE;
if cond { z = x; } else { z = y; }
fputc('a', z); fclose(x); pclose(y);
// After transformation
let x: BufWriter<File> = BufWriter::new(File(..));
let y: Child = Command::new(..).spawn();
let z: *mut dyn Write;
if cond { z = &raw mut x; } else { z = &raw mut y; }
z.write_all(&['a']); drop(x); drop(y);
\end{lstlisting}

\noindent
%
Here, \C{x} and \C{y} have base types because each has a single origin and is
closed, while \C{z} is a \C{dyn} pointer because it has multiple origins and is
not closed, necessitating proper conversions in the assignments from \C{x} and
\C{y} to \C{z}.

\subsection{Propagating Errors}
\label{sec:transformation:errors}

In this step, we transform the program to propagate errors.
%
We first define a local variable to store a stream indicator in each function
that either checks the indicator or contains an error source.
%
Each variable is initialized to \C{0}, indicating no error.
%
We also replace each \C{feof}/\C{ferror} call with an expression that reads this
error variable.

We then transform each Rust std API call acting as a source to check whether the
returned \C{Result} is \C{Err}.
%
If so, we update the error variable to a non-zero value, indicating an error.

Finally, we transform user-defined functions to propagate errors.
%
If a source is in a callee, we modify the callee to return a tuple consisting of
the original return value and the error value;
%
at the call site, we use the returned error value to update the caller's error
variable.
%
Conversely, if a source is in a caller, we transform the callee to take an
additional parameter, and modify the caller to pass the error variable as an
argument.

Below is an example showing the result of the transformation, where \C{bar}
contains a source:

\begin{lstlisting}[language=Rust]
// Before transformation
fn bar(x:*mut FILE)->i32 {fputc('a', x); return ..;}
bar(y); if ferror(y) != 0 { .. }
// After transformation
fn bar(x: impl Write) -> (i32, i32) { let e = 0;
    if x.write_all(&['a']).is_err() { e = 1; }
    return (.., e); }
let e=0; let (_, _e)=bar(&mut y); e=_e; if e!=0 {..}
\end{lstlisting}

Unfortunately, the transformed code does not follow Rust's idioms for error
handling.
%although it preserves the behavior of the original program.
%
The idiomatic approach would be to transform functions to take or return
\C{Option}/\C{Result} values instead of integers.
%
Since this work focuses on preserving semantics during I/O translation, we leave
improving idiomaticity to future work.
